%==============================================================================
% APPENDIX B — USER STUDY QUESTIONNAIRE
%==============================================================================
\chapter{User Study Instruments}
\label{app:user-study}

\section*{SUS Questionnaire (System Usability Scale)}

Participants rated the following 10 statements on a 5-point Likert scale (1 = Strongly Disagree, 5 = Strongly Agree):

\begin{enumerate}
  \item I think that I would like to use this system frequently.
  \item I found the system unnecessarily complex.
  \item I thought the system was easy to use.
  \item I think that I would need the support of a technical person to be able to use this system.
  \item I found the various functions in this system were well integrated.
  \item I thought there was too much inconsistency in this system.
  \item I would imagine that most people would learn to use this system very quickly.
  \item I found the system very cumbersome to use.
  \item I felt very confident using the system.
  \item I needed to learn a lot of things before I could get going with this system.
\end{enumerate}

\section*{TAM Questionnaire (Technology Acceptance Model)}

\textbf{Perceived Usefulness (PU)} --- 5 items, 7-point Likert (1 = Extremely Unlikely, 7 = Extremely Likely):
\begin{enumerate}
  \item Using AntiGravity improves my performance in finding information.
  \item Using AntiGravity increases my productivity.
  \item Using AntiGravity enhances my effectiveness on the job.
  \item AntiGravity is useful in my job.
  \item Using AntiGravity saves me time when searching for information.
\end{enumerate}

\textbf{Perceived Ease of Use (PEOU)} --- 4 items, 7-point Likert:
\begin{enumerate}
  \item Learning to operate AntiGravity is easy for me.
  \item I find it easy to get AntiGravity to do what I want it to do.
  \item My interaction with AntiGravity is clear and understandable.
  \item I find AntiGravity easy to use.
\end{enumerate}

\section*{Task Scenarios}

The five task scenarios used in the task-based evaluation:

\begin{enumerate}
  \item \textbf{T1 (Factoid):} ``Find the retention period for project closure reports in the DNSI document management policy.''
  \item \textbf{T2 (Multi-hop):} ``What are the security requirements for deploying a new application that processes personal data on the CERP infrastructure?''
  \item \textbf{T3 (Procedural):} ``Find the step-by-step procedure for requesting access to a new SharePoint site.''
  \item \textbf{T4 (Multi-hop):} ``Compare the roles and responsibilities of the Project Manager and the Technical Lead in the DNSI project governance framework.''
  \item \textbf{T5 (Unanswerable):} ``What is the direct phone number of the DNSI's external cybersecurity auditor?''
\end{enumerate}

\section*{Semi-Structured Interview Guide}

\begin{enumerate}
  \item How would you describe your overall experience using AntiGravity?
  \item For which types of tasks would you prefer to use AntiGravity over SharePoint search?
  \item What features did you find most valuable? Why?
  \item What aspects of the system frustrated you?
  \item How much do you trust the responses provided by AntiGravity? What factors influence your trust?
  \item Would you use AntiGravity in your daily work if it were permanently deployed? Under what conditions?
\end{enumerate}
